\section*{Erd\H{o}s Problem \#273}

\subsection*{1) ROUND-2 OBJECTIVE}

Path (C): isolate the exact defect in the statement used in Round-1, prove the literal page sentence completely, and identify the minimal correction needed to recover the genuinely nontrivial problem.

The reason this is the right direction is that Round-1 already showed that the distinct-modulus version is nontrivial and presently unresolved, whereas the literal sentence on the problem page omits the distinctness convention.

\subsection*{2) Round-1 FOUNDATION USED}

I use the following Round-1 results.

\begin{enumerate}
\item Round-1 observed that if repeated moduli are allowed, then the modulus $4=5-1$ already gives a trivial cover.
\item Round-1 proved for the corrected distinct-modulus version that any such cover must have at least $18$ congruences and at least one modulus $\ge 70$.
\end{enumerate}

Only (2) is carried forward as information about the corrected problem.

\subsection*{3) NEW INSIGHT / TOOL (ROUND-2)}

The new point is that the ``distinct moduli'' correction is not ad hoc.  In one of Erd\H{o}s's own expository sources, a covering system is explicitly written in the form
\[
a_i \ (\bmod n_i), \qquad n_1<n_2<\cdots<n_k.
\]
So the historical convention in that source already bakes in distinctness of the moduli.

Therefore there are really two different statements:

\begin{enumerate}
\item the \emph{literal standalone sentence} on the webpage, which does not say ``distinct'', and
\item the \emph{historically intended} covering-system question under Erd\H{o}s's convention, where the moduli are distinct.
\end{enumerate}

\subsection*{4) ATTACK PLAN (ROUND-2)}

The exact Round-1 gap was this: it corrected the problem to the distinct-modulus version, but did not yet settle whether that correction was logically forced or merely stylistic.

To close that gap I will:

\begin{enumerate}
\item prove the literal sentence directly by an explicit cover;
\item justify the minimal corrected statement using Erd\H{o}s's own distinct-modulus convention;
\item record that the corrected statement is exactly the open version to which the Round-1 lower bounds apply.
\end{enumerate}

\subsection*{5) WORK (ROUND-2)}

\paragraph{Literal page sentence: complete affirmative solution.}
Consider the four congruences
\[
x\equiv 0 \ (\bmod 4),\qquad
x\equiv 1 \ (\bmod 4),\qquad
x\equiv 2 \ (\bmod 4),\qquad
x\equiv 3 \ (\bmod 4).
\]
Every integer has one of the four residues modulo $4$, so these four congruences cover all integers.

Their common modulus is $4$, and
\[
4=5-1
\]
with $5$ prime and $5\ge 5$.  Hence every modulus is of the required form $p-1$ for a prime $p\ge 5$.

So the literal sentence

\medskip
\centerline{\emph{``Is there a covering system all of whose moduli are of the form $p-1$ for some primes $p\ge 5$?''}}
\medskip

has the answer \emph{Yes}, by the explicit cover above.

\paragraph{Minimal correction.}
Under Erd\H{o}s's own distinct-modulus convention, the intended nontrivial statement is:

\medskip
\centerline{\emph{Does there exist a covering system with pairwise distinct moduli, each of the form $p-1$ for a prime $p\ge 5$?}}
\medskip

This is exactly the corrected statement adopted in Round-1.

\paragraph{Status of the corrected statement.}
Round-1 already proved two necessary conditions for that corrected problem:

\begin{enumerate}
\item at least $18$ congruences are needed;
\item at least one modulus must be $\ge 70$.
\end{enumerate}

No construction or impossibility proof is presently known to me beyond those Round-1 obstructions.

\subsection*{6) ADVERSARIAL VERIFICATION}

There are two possible objections.

\begin{enumerate}
\item \emph{``A covering system already means distinct moduli.''}  
If one adopts that convention from the start, then the four congruences modulo $4$ are indeed irrelevant to the intended problem.  But that only confirms the need to state the convention explicitly.  The standalone webpage sentence does not do so.
\item \emph{``Then the problem has not really been solved.''}  
Correct: the \emph{corrected distinct-modulus problem} remains open.  What \emph{has} been settled in this round is the logical status of the sentence as written: taken literally, it is trivial; taken in Erd\H{o}s's convention, it becomes the open distinct-modulus problem studied in Round-1.
\end{enumerate}

So no contradiction remains: the defect is a missing convention, and the minimal repair is exactly the distinctness condition.

\subsection*{7) FINAL}

\textbf{UNRESOLVED (BUT STRICTLY ADVANCED).}

The literal statement is completely solved in the affirmative by the modulus-$4$ example.  The historically intended distinct-modulus version is the minimal corrected statement, and that corrected version remains open.  Round-1's lower bounds continue to hold for it.

\subsection*{8) COMPLETION ESTIMATE (MANDATORY)}

\textbf{COMPLETION: 65\%}

\subsection*{9) REFERENCES}

\begin{enumerate}
\item P. Erd\H{o}s, \emph{Some of my favourite problems in number theory}.  The covering-systems section explicitly writes the moduli in strictly increasing order.
\item T. F. Bloom, \emph{Erd\H{o}s Problem \#273}.  Current webpage statement and current status check.
\item Round-1 output for Problem \#273 (reciprocal-sum obstruction and the lower bounds ``at least $18$ moduli'' and ``some modulus $\ge 70$'').
\end{enumerate}



